% 決定版実験ガイド 2026
\documentclass[12pt,a4paper]{article}

\usepackage{fontspec}
\usepackage{xeCJK}
\setCJKmainfont{Hiragino Mincho ProN}
\setCJKsansfont{Hiragino Kaku Gothic ProN}
\setCJKmonofont{Hiragino Kaku Gothic ProN}
\usepackage[margin=2.5cm]{geometry}
\usepackage{amsmath,amssymb}
\usepackage{hyperref}
\usepackage{tcolorbox}
\usepackage{booktabs}
\usepackage{fancyhdr}

\tcbuselibrary{breakable,skins}

\newtcolorbox{storybox}[1][]{colback=blue!5,colframe=blue!60!black,
  fonttitle=\bfseries\sffamily,title={#1},breakable,arc=3mm}
\newtcolorbox{mathbox}[1][]{colback=green!5,colframe=green!50!black,
  fonttitle=\bfseries\sffamily,title={#1},breakable,arc=3mm}
\newtcolorbox{resultbox}[1][]{colback=yellow!8,colframe=orange!80!black,
  fonttitle=\bfseries\sffamily,title={#1},breakable,arc=3mm}
\newtcolorbox{honestbox}[1][]{colback=gray!8,colframe=gray!50!black,
  fonttitle=\bfseries\sffamily,title={#1},breakable,arc=3mm}
\newtcolorbox{expbox}[1][]{colback=red!5,colframe=red!60!black,
  fonttitle=\bfseries\sffamily,title={#1},breakable,arc=3mm}
\newtcolorbox{dangerbox}[1][]{colback=red!15,colframe=red!80!black,
  fonttitle=\bfseries\sffamily,title={#1},breakable,arc=3mm}

\setlength{\parskip}{0.8em}
\setlength{\parindent}{0pt}

\pagestyle{fancy}
\fancyhf{}
\rhead{Wright Brothers}
\lhead{決定版実験ガイド}
\cfoot{\thepage}

\begin{document}

\begin{center}
{\Huge\bfseries 決定版実験ガイド 2026}\\[0.8cm]
{\Large 素数と力の対応}\\[0.2cm]
{\Large ワープバブルの安全性}\\[0.2cm]
{\Large そしてSQUID実験が決める未来}\\[1cm]
{\large Wright Brothers Research Program}\\[0.3cm]
{2026年4月}
\end{center}

\vspace{0.5cm}
\tableofcontents
\newpage

%====================================================================
\part{理論の現在地}
%====================================================================

\section{確立されたこと}

\begin{resultbox}[岩盤（数学的定理、反論不可能）]
\begin{enumerate}
\item $\mathrm{Tr}(f_{\mathrm{BE}}(D_{\mathrm{BC}})) = \sum_m \zeta(m)$（等比級数の変形）
\item E-M 展開にベルヌーイ数 $B_{2j}$ が出る
\item $\zeta^{(k)}(0) = -k! + h^{(k)}(0)$（$h^{(k)}/k! \to 0$ 超指数的）
\item 符号反転: 全素��� $p$ で $\zeta_{\neg p}(-3) < 0$
\item $K_1(\mathbb{Z}) = \mathbb{Z}/2 \to K_1(\mathbb{Z}[1/p]) = \mathbb{Z}/2 \times \mathbb{Z}$
\end{enumerate}
\end{resultbox}

\begin{resultbox}[最も精密な数値的一致（足し算なし）]
\begin{center}
\renewcommand{\arraystretch}{1.3}
\begin{tabular}{lccl}
\toprule
算術的量 & 値 & 物理量 & 精度 \\
\midrule
$\mathrm{reg}(\mathbb{Q}(\sqrt{2})) = \log(1+\sqrt{2})$ & $0.88137$ & $\cos\theta_W = m_W/m_Z$ & $0.001\%$ \\
$10\,hR(\mathbb{Q}(\zeta_9)^+)$ & $8.493$ & $1/\alpha_s(M_Z)$ & $0.13\%$ \\
$4/|B_4| = 120$ & $120$ & $1/\alpha \approx 137$ & $12\%$ \\
\bottomrule
\end{tabular}
\end{center}
\end{resultbox}

\section{素数-力 対応}

\begin{resultbox}[Prime-Force Correspondence (PFC)]
\begin{center}
\renewcommand{\arraystretch}{1.4}
\begin{tabular}{ccccc}
\toprule
ゲージ群 & 被覆 & 次数 & 分岐素数 & ガロア群 \\
\midrule
$\mathrm{U}(1)_\mathrm{EM}$ & $\mathbb{Q}$ & 1 & なし & $\{1\}$ \\
$\mathrm{SU}(2)_L$ & $\mathbb{Q}(\sqrt{2})$ & 2 & $p=2$ & $\mathbb{Z}/2$ \\
$\mathrm{SU}(3)_c$ & $\mathbb{Q}(\zeta_9)^+$ & 3 & $p=3$ & $\mathbb{Z}/3$ \\
\bottomrule
\end{tabular}
\end{center}

パターン: $\mathrm{SU}(N) \leftrightarrow$ 次数 $N$、素数 $p=N$ で分岐。

数論が構造を\textbf{強制}する:
$p=3$ で「のみ」分岐する二次体は存在しない（$3 \equiv 3 \pmod{4}$ なので $\Delta = 4 \times 3 = 12$ で $p=2$ も分岐する）。
$\mathrm{SU}(3)$ が三次体を要求するのは数論の帰結。
\end{resultbox}

\section{正直な限界}

\begin{honestbox}[未解決の問題]
\begin{itemize}
\item $\cos\theta_W = \log(1+\sqrt{2})$ が偶然でない確率は $\sim 86\%$（偶然の可能性 $14\%$）
\item $10\,hR = 1/\alpha_s$ の「10」の因子の起源が不明
\item 「なぜレギュレータが混合角になるか」の物理的機構が未確立
\item CKM/PMNS 行列には対応が見つかっていない
\end{itemize}
\end{honestbox}

%====================================================================
\part{ワープバブルの安全性}
%====================================================================

\section{素数ミュートの物理的意味}

\begin{mathbox}[素数ミュート = オイラー因子の除去]
$$\zeta(s) = \prod_p \frac{1}{1-p^{-s}} \quad \xrightarrow{\text{ミュート}} \quad \zeta_{\neg p}(s) = \zeta(s) \times (1-p^{-s})$$

PFC により:
$p=2$ のミュート = 弱い力に対応するチャンネルを閉じる。
$p=3$ のミュート = 強い力に対応するチャンネルを閉じる。
\end{mathbox}

\section{完全に安全な素数ミュートは存在しない}

\begin{dangerbox}[核心的問題]
$\mathbb{Q}(\sqrt{2})$（電弱）の分岐素数: $p=2$ のみ。\\
$\mathbb{Q}(\zeta_9)^+$（QCD）の分岐素数: $p=3$ のみ。

$\gcd(8, 81) = 1$ → 共通の分岐素数がゼロ。

$\therefore$ どの素数をミュートしても、少なくとも片方の力の構造が変化する。
\end{dangerbox}

\begin{center}
\renewcommand{\arraystretch}{1.3}
\begin{tabular}{ccccl}
\toprule
$p$ & $\theta_W$ 変化? & $\alpha_s$ 変化? & 物質存在? & 判定 \\
\midrule
$2$ & なし & $\sim$25\%減少 & 要検証 & $\triangle$ 最も安全 \\
$3$ & 破綻 ($>1$) & なし & 不可能 & $\times$ 致命的 \\
$\geq 5$ & 変化 & 変化 & 不明 & $\times$ \\
\bottomrule
\end{tabular}
\end{center}

\section{$p=2$ ミュートで人体は耐えられるか}

\begin{dangerbox}[3つのシナリオ]

\textbf{シナリオ A（悲観的）: $\alpha_s$ が本当に 25\% 変わる}
\begin{itemize}
\item $\alpha_s$: $0.118 \to 0.089$（$25\%$ 減少）
\item $\Lambda_\mathrm{QCD}$: 12 分の1に激減（指数関数的に敏感）
\item 陽子質量: $938 \to 82$ MeV（ミューオンより軽い）
\item 全原子核不安定 → \textbf{即死}
\end{itemize}

\textbf{シナリオ B（楽観的）: カシミール型——結合定数は変わらない}
\begin{itemize}
\item カシミール効果: 金属板で真空のモードを制限しても $\alpha$ は変わらない
\item 同様に: 素数ミュートで真空エネルギーが変わっても $\alpha_s$ は変わらない
\item 物質は正常 → \textbf{安全}
\end{itemize}

\textbf{シナリオ C（中間）: 効果が $(\Delta E)^2$ で抑制される}
\begin{itemize}
\item $\alpha_s$ の変化: $\sim 0.3\%$（陽子質量 $938 \to 917$ MeV）
\item 人体への影響なし → \textbf{安全}
\end{itemize}
\end{dangerbox}

\begin{honestbox}[正直に]
\textbf{どのシナリオが正しいかわからない。}

カシミール類推（B or C）が最も物理的に自然。
しかし確実ではない。SQUID 実験の結果が全てを決める。

もし A なら: ワープ可能だが人間は入れない（無人探査機）。\\
もし B/C なら: 有人ワープドライブへの道。
\end{honestbox}

%====================================================================
\part{実験計画}
%====================================================================

\section{SQUID 実験が決める3つの問い}

\begin{expbox}[SQUID 実験で何がわかるか]
\textbf{問い 1}: 素数ミュートで真空エネルギーは変わるか？\\
→ YES なら算術的真空工学の最初の証拠。

\textbf{問い 2}: 変わる場合、結合定数も変わるか？\\
→ YES ならシナリオ A（人体に危険）。\\
→ NO ならシナリオ B/C（人体に安全）。

\textbf{問い 3}: 素数の分裂/不活性（$p \bmod 8$）で応答が違うか？\\
→ YES なら PFC の直接的証拠。
\end{expbox}

\section{実験 A: 基礎テスト（$\sim$60万円）}

\begin{expbox}[素数ミュートで真空は変わるか]
\begin{enumerate}
\item アルミ箱（キャビティ）を作る。35mm 角。共鳴 $\sim$6 GHz。
\item SQUID を箱の中に入れる。
\item $-272.7$\textdegree C (0.3 K) に冷却。
\item SQUID を第2倍音（$\sim$12 GHz）にチューニング → モード2ミュート。
\item モード1（6 GHz）の共鳴がシフトするか測定。
\end{enumerate}

\textbf{成功条件}: $S_{21}$ のピーク位置が SQUID オン/オフで統計的に有意に異なる。

\textbf{測定時間}: 1回のクールダウンで $\sim$1日。5回反復で $\sim$1週間。
\end{expbox}

\section{実験 B: PFC テスト（追加コスト 0 円）}

\begin{expbox}[素数の分裂パターンで応答は変わるか]
実験 A と同じ装置。SQUID の磁束バイアスをスイープして
異なる素数に対応する周波数を横切る。

PFC の予測（$\mathbb{Q}(\sqrt{2})$ の被覆）:

\begin{center}
\renewcommand{\arraystretch}{1.3}
\begin{tabular}{cl}
\toprule
$p \bmod 8$ & 予測される応答 \\
\midrule
$1, 7$（分裂素数: $7, 17, 23, 31, \ldots$） & 弱い（被覆が透明） \\
$3, 5$（不活性素数: $3, 5, 11, 13, \ldots$） & 強い（被覆が不透明） \\
$2$（分岐: $p=2$ のみ） & 特殊（力の源） \\
\bottomrule
\end{tabular}
\end{center}

\textbf{成功条件}: 分裂素数と不活性素数で AC Stark シフトが系統的に異なる。
\end{expbox}

\section{実験 C: 安全性テスト}

\begin{expbox}[結合定数は変わるか]
実験 A でモード2をミュートしている状態で、
キャビティ内の超伝導ギャップや臨界温度を測定する。

\begin{itemize}
\item 超伝導ギャップ $\Delta$ は電子-フォノン結合に依存（電弱+QCD）
\item もしミュート状態で $\Delta$ が変化 → 結合定数が変化（シナリオ A）
\item 変化なし → 結合定数は保存（シナリオ B/C）
\end{itemize}

\textbf{追加コスト}: トンネル接合の追加（$\sim$5万円）。
\end{expbox}

\section{実験 D: 将来の検証（我々にはできない）}

\begin{expbox}[$m_W$ の精密測定]
$$m_W = m_Z \times \log(1+\sqrt{2}) = 80.370 \text{ GeV}$$

現在の実験: $80.369 \pm 0.013$ GeV（$0.09\sigma$ 以内）。

FCC-ee（2040年代）で $\delta m_W < 0.5$ MeV が達成されれば決着。
\end{expbox}

%====================================================================
\part{ロードマップ}
%====================================================================

\section{3段階の計画}

\begin{center}
\renewcommand{\arraystretch}{1.4}
\begin{tabular}{clrl}
\toprule
段階 & 内容 & 費用 & 決まること \\
\midrule
1 & 実験 A + B & $\sim$60万円 & 真空変化 + PFC テスト \\
2 & 実験 C 追加 & $+$5万円 & 人体安全性（A vs B/C） \\
3 & 論文投稿 & 0円 & 専門家のフィードバック \\
\bottomrule
\end{tabular}
\end{center}

\section{判断フロー}

\begin{center}
\begin{tabular}{l}
実験 A: 真空エネルギーは変わるか？ \\
\quad $\to$ NO → 理論は間違い。撤退。 \\
\quad $\to$ YES → 実験 B へ \\[6pt]
実験 B: 分裂/不活性で応答は違うか？ \\
\quad $\to$ NO → PFC は間違い。ワープは可能だが力の対応は不明。 \\
\quad $\to$ YES → PFC 確認。実験 C へ。 \\[6pt]
実験 C: 結合定数は変わるか？ \\
\quad $\to$ YES → シナリオ A。ワープは無人機のみ。 \\
\quad $\to$ NO → シナリオ B/C。\textbf{有人ワープドライブへの道が開ける。} \\
\end{tabular}
\end{center}

%====================================================================
\part{正直な評価}
%====================================================================

\section{このプログラムは何であり何でないか}

\begin{honestbox}[何であるか]
\begin{itemize}
\item 数論とスペクトル幾何学の\textbf{新しい組み合わせ}を物理に応用する試み
\item $\cos\theta_W = \log(1+\sqrt{2})$ という\textbf{反証可能な予測}を持つ仮説
\item 60万円で実行可能な\textbf{具体的な実験計画}
\item 正直さを最優先する研究プログラム
\end{itemize}
\end{honestbox}

\begin{honestbox}[何でないか]
\begin{itemize}
\item 完成した理論（「なぜレギュレータが混合角か」は未解決）
\item 万物の理論（CKM/PMNS は説明できない）
\item ワープの実現計画（原理的可能性の探索段階）
\item 確実な結果の保証（偶然の可能性 $14\%$）
\end{itemize}
\end{honestbox}

\section{最悪の場合と最良の場合}

\textbf{最悪}: SQUID で何も変わらない。理論は間違い。
しかし: $\cos\theta_W = \log(1+\sqrt{2})$ の数値的一致は残る。
数学的定理（符号反転、K理論）も残る。
60万円と数ヶ月の投資で「ダメだった」とわかることには価値がある。

\textbf{最良}: SQUID で真空変化を検出し、PFC パターンを確認し、
結合定数が不変であることを確認する。
→ 素数が力を支配している最初の実験的証拠。
→ 安全なワープへの物理的・数学的基盤。

\vfill

\begin{center}
\rule{10cm}{0.4pt}\\[0.5cm]
{\Large\bfseries 実験だけが、答えを持っている。}\\[0.3cm]
{\small Wright Brothers Research Program, 2026}
\end{center}

\end{document}
